\section{Assembling the recursion}\label{sec:recursion}

The two field encodings now reduce synthesis of $V$ to synthesis of the
smaller controlled unitary $F(z)$.  The reduction uses the smaller circuit
exactly once.  This fact keeps the total circuit size polynomial even though
the number of recursion levels grows with $n$.

\begin{lemma}[One synthesis step]\label{rec:step}
Let $V$ be a unitary on $\C^{k+s}$ in the block form of
Proposition~\ref{contract:block}, and let $\norm D\leq r<1$.
Use the fields and encoders of Sections~\ref{sec:fields} and~\ref{sec:encoders}
with truncation length $L$.  Suppose a circuit synthesizes the family of
$k$-dimensional unitaries $F(z)$, controlled by $z$, with clean-workspace
operator norm error $\delta'$.  Then one use of that circuit, preceded by
$\mathsf U_E$ and followed by $\mathsf U_R^*$, synthesizes $V$ with
clean-workspace error at most
\begin{equation}\label{rec:step-error}
 \delta'+\frac{2r^L}{1-r}+10r^{2L}.
\end{equation}
The assertion also holds coherently for a family $V_\lambda$, with the
smaller circuit controlled by both $\lambda$ and $z$.
\end{lemma}

\begin{proof}
Let $\iota$ embed the current index space into the full register layout by
initializing all working qubits to zero.  Include the fresh working qubits
of the smaller circuit in this definition.  Let $\jmath$ embed the field space
$\C^{\mathcal G}\otimes\C^k$ into the encoder's good subspace, again with
the smaller circuit's fresh workspace zero.  Write $\mathsf C$ for the
smaller synthesis circuit acting on its target, control, and fresh working
registers, and as the identity on the remaining registers.  Its approximation
guarantee implies
\begin{equation}\label{rec:restricted-error}
 \norm{\mathsf C\jmath-\jmath\widetilde F}\leq\delta'.
\end{equation}
Indeed, the good subspace has a valid port index, and $\widetilde F$ applies
$F(z)$ with the phase register as control.

Put $\tau=r^L/(1-r)$.  Proposition~\ref{encode:implementation} gives
\[
 \norm{\mathsf U_E\iota-\jmath\mathcal E_L}\leq5r^{2L},
 \qquad
 \norm{\mathsf U_R\iota-\jmath\mathcal R_L}\leq5r^{2L}.
\]
The field maps are contractions by Proposition~\ref{fields:gram}, and their
intertwining bound is~\eqref{fields:intertwine}.  Consequently,
\begin{align*}
 &\norm{\mathsf C\mathsf U_E\iota-\mathsf U_R\iota V}\\
 &\quad\leq
 \norm{\mathsf C(\mathsf U_E\iota-\jmath\mathcal E_L)}
 +\norm{(\mathsf C\jmath-\jmath\widetilde F)\mathcal E_L}\\
 &\qquad\quad
 +\norm{\jmath(\widetilde F\mathcal E_L-\mathcal R_LV)}
 +\norm{(\jmath\mathcal R_L-\mathsf U_R\iota)V}\\
 &\quad\leq5r^{2L}+\delta'+2\tau+5r^{2L}.
\end{align*}
Multiplication by $\mathsf U_R^*$ proves~\eqref{rec:step-error}.
This calculation invokes the smaller circuit's approximation only on the
ideal field-map output, whose norm is at most one.  Any component of the
actual encoded state outside the good subspace is covered by the first
error term and is not enlarged by $\mathsf C$.

For additional labels $\lambda$, all maps are block diagonal in the label
basis and all bounds are uniform over the labels.  Their operator norms on
the full controlled space are the maxima of the block norms.  The same
calculation therefore applies to arbitrary coherent superpositions of labels.
\end{proof}

Figure~\ref{fig:recursion} shows the comparison made in this proof.  The
amplification inside $\mathsf U_E$ and $\mathsf U_R$ repeats only their own
state-preparation circuits.  It never repeats $\mathsf C$.
\begin{figure}[tb]
\centering
\begin{tikzpicture}[>=Latex, font=\small,
  box/.style={draw,rounded corners=2pt,minimum height=10mm,align=center},
  state/.style={align=center,font=\footnotesize}]
  \node[box,minimum width=19mm] (e) at (0,0) {$\mathsf U_E$};
  \node[box,minimum width=39mm] (f) at (4.8,0)
       {$\mathsf C$\,: one recursive call};
  \node[box,minimum width=19mm] (r) at (9.6,0) {$\mathsf U_R^*$};
  \draw[->] (-1.8,0) -- (e.west);
  \draw[->] (e.east) -- (f.west);
  \draw[->] (f.east) -- (r.west);
  \draw[->] (r.east) -- (11.4,0);
  \node[state] at (-1.45,.58) {$\iota x$};
  \node[state] at (2.05,.58) {$\jmath\mathcal E_Lx$};
  \node[state] at (7.55,.58) {$\jmath\mathcal R_LVx$};
  \node[state] at (11.02,.58) {$\iota Vx$};
  \node[state] at (0,-.92) {local encoding};
  \node[state] at (4.8,-.92) {order $k<d$; controlled by $z$\\
                           fresh workspace starts at zero};
  \node[state] at (9.6,-.92) {local decoding};
  \draw[decorate,decoration={brace,mirror,amplitude=4pt}]
       (1.22,-1.7) -- (8.38,-1.7);
  \node[state] at (4.8,-2.12)
       {comparison states lie in the fixed good subspace};
\end{tikzpicture}
\caption{One recursion level.  Labels above the arrows are the comparison
states used in Lemma~\ref{rec:step}; the intermediate states are approximate,
as quantified there.  The phase register supplies the control $z$ to the
single smaller synthesis circuit.  The local encoders contain their own
amplification rounds, while the smaller circuit occurs only once.  The map
$\iota$ initializes all work registers to zero, so the final comparison
also records their cleanup.}
\label{fig:recursion}
\end{figure}


We can now choose one set of parameters for all recursion levels.  Keep
$M,r,J$ from~\eqref{contract:parameters} and~\eqref{contract:J}, and set
\begin{equation}\label{rec:parameters}
 L=10000J,\qquad
 b=nL+\lceil\log_2(L+1)\rceil,\qquad
 K=b(L+1).
\end{equation}
Thus $K$ is the number of qubits in each field's phase register.

\begin{proposition}[Uniform programmable synthesis]\label{rec:skeleton}
Let $n\geq1$, $1\leq d\leq2^n$, and $t\geq0$.  There is a programmable-slot
circuit structure, determined by $n,d,t$, with the following property.
For every family $(U_\lambda)_{\lambda\in\{0,1\}^t}$ of unitaries on
$\C^d$, its slot tables can be chosen so that the circuit implements the
controlled unitary $U_\lambda$, preserves the labels, and has clean-workspace
operator norm error less than $1/100$ in the sense of~\eqref{slot:clean-error}.
The circuit uses
\begin{equation}\label{rec:resources}
 t+O((n+1)^{10})\quad\text{qubits},
 \qquad O((n+1)^{13})\quad\text{slots}.
\end{equation}
Its structure, including each slot's ordered control list, is generated
deterministically in time polynomial in $n+t$ from $n,d,t$.
\end{proposition}

\begin{proof}
We recurse on the valid order $d$, retaining $n$-bit index and port registers
at every level.  If $d=1$, each $U_\lambda$ is a scalar phase, implemented
exactly by one slot on an index bit that is zero throughout the valid input
subspace.

For $d\geq2$, Proposition~\ref{contract:block} gives
$V_\lambda=P_{\mathrm L,\lambda}U_\lambda P_{\mathrm R,\lambda}$ with
internal order $s$ and port order $k=d-s$.  Apply the one-step construction
of Lemma~\ref{rec:step}, prescribing the recursive family to be
$F_\lambda(z)$ and adding $z$ to the existing control labels.  Supply fresh
working qubits for this recursive circuit.  To implement $U_\lambda$, apply
$P_{\mathrm R,\lambda}^*$ before this construction and
$P_{\mathrm L,\lambda}^*$ afterward.  These transformations have exact slot
implementations, so they introduce no further error.  Their order is correct
because $P_{\mathrm L,\lambda}^*V_\lambda P_{\mathrm R,\lambda}^*=U_\lambda$.

At each step, Lemma~\ref{rec:step} adds at most
\[
 \frac{2r^L}{1-r}+10r^{2L}
 =8r^L+10r^{2L}\leq18r^L
\]
to the error of the smaller circuit.  There are fewer than $J$ steps by
Lemma~\ref{contract:depth}.  Bernoulli's inequality gives
$(4/3)^L\geq1+L/3$, and hence $r^L\leq3/L$.  The total error is therefore
\begin{equation}\label{rec:total-error}
 18Jr^L\leq\frac{54J}{L}=\frac{54}{10000}<\frac1{100}.
\end{equation}
In particular, all working qubits, including those introduced at inner
recursion levels, return to zero in the ideal comparison map.

We next count the circuit resources.  Write $N=n+1$ and
$h=\lceil\log_2(L+1)\rceil$.  Our parameters satisfy
\[
 J,L=O(N^3),\qquad b=O(N^4),\qquad K=O(N^7),\qquad h=O(N).
\]
Proposition~\ref{encode:implementation} uses $O(K+n+h)$ fresh qubits for a
pair of encoders and $O(L(K+n+h)+K)$ slots for each encoder.  The two encoders
at one level share the same register layout.  The pre- and post-conditioning
operations require only $O(n)$ more clean working qubits and slots.
All registers from an outer level remain allocated during an inner call,
so summing over the at most $J$ levels gives
\[
 t+O(J(K+n+h))=t+O(N^{10})
\]
total qubits.  Summing the slots gives
\[
 O\bigl(J[L(K+n+h)+K+n]\bigr)=O(N^{13}).
\]
These sums suffice because there is only one recursive call at a level.
The local amplification repetitions have already been included in the
factor $L$ in the encoder cost.

Finally, the structure can be generated without computing any of its table
contents.  The successive orders $d,q,s,k$ have $O(n)$-bit descriptions and
are computed in polynomial time.  Register allocations, state-preparation
bit order, Fourier transforms, reflections, and the fixed number $L$ of
amplification rounds depend only on these sizes.  Every operation depending
on $U_\lambda$---the signs, permutations, state-preparation amplitudes,
support lists, and recursively prescribed matrices---changes only a slot
table.  At most $t+O(N^{10})$ register positions occur in any ordered control
list, so writing all such lists for $O(N^{13})$ slots still takes polynomial
time in $n+t$.  No list of possible label values is generated.  This proves
both the uniformity assertion and~\eqref{rec:resources}.
\end{proof}
